\documentclass[11pt]{article}
\usepackage[a4paper,margin=1in]{geometry}
\usepackage{amsmath,amssymb,amsthm,mathtools}
\usepackage[hidelinks]{hyperref}
\usepackage{booktabs}
\usepackage{array}
\pdftrailerid{}

\newtheorem{theorem}{Theorem}
\newtheorem{proposition}[theorem]{Proposition}
\newtheorem{lemma}[theorem]{Lemma}
\newtheorem{corollary}[theorem]{Corollary}
\newtheorem{remark}[theorem]{Remark}
\newcommand{\Q}{\mathbb Q}
\newcommand{\Z}{\mathbb Z}
\newcommand{\Cl}{\operatorname{Cl}}
\newcommand{\Gal}{\operatorname{Gal}}
\newcommand{\Norm}{\operatorname{N}}
\newcommand{\ord}{\operatorname{ord}}
\newcommand{\ab}{\mathrm{ab}}

\title{Effective Archimedean Stark Theorems over Real Quadratic Fields:\\
Quadratic Support, Shintani Transfer, and Higher-Order Packets}
\author{Hainan Zhao\\
\small Independent Researcher\\
\small\texttt{hainzhao@gmail.com}}
\date{30 July 2026}

\begin{document}
\maketitle

\begin{abstract}
We prove a uniform formula for every one-place real-quadratic Stark
invariant whose differenced Fourier support is quadratic.  The formula
is an explicit product of relative units, with rational exponents
computed from class numbers, imprimitive Euler factors, and an exact
relative-regulator index.  Beyond this quadratic floor, we translate
Shintani's index-two theorem into the one-place notation, give a
general all-embeddings height-rigidity criterion, and use them to
identify seven explicit packets with support orders up to ten.  The
examples include two order-six packets, an order-six packet at a
ramified \(3\)-power conductor, and an order-ten packet.  Two
structural lemmas prove that an absolutely abelian one-place ray field
has trivial differenced support and that nontrivial support forces
even Shintani index.  The latter is independently consistent with all
446 genuinely reconstructed odd-index cases in a separate census.
Exact and Arb certificates accompany the written proofs; no unproved
Stark conjecture or floating-point recognition enters a proof.
\end{abstract}

\section{Introduction and claim boundary}

Let $K$ be real quadratic and let
\(\mathfrak m=\mathfrak f\infty_2\) be a ray modulus containing one
real place.  For \(K=\Q(\sqrt D)\) with \(D>0\), our convention is
that \(\infty_2\) is the negative-square-root embedding
\(\sqrt D\mapsto-\sqrt D\), which is the first real place in the
pinned PARI ordering; \(\infty_1\) is the positive-square-root
embedding.  For \(A\in\Cl_{\mathfrak m}(K)\), let \(R\) be the sign
class and put
\[
 Z_{\mathfrak m}(s,A)
 =\zeta_{\mathfrak m}(s,A)
  -\zeta_{\mathfrak m}(s,RA),\qquad
 X_A=\exp Z'_{\mathfrak m}(0,A).
\tag{1}
\]
The numbers \(X_A\) are positive real analytic invariants.  We call
the finite Artin-indexed tuple
\[
 \mathcal X_{\mathfrak m}=(X_A)_{A\in\Cl_{\mathfrak m}(K)}
\]
the \emph{one-place packet}; repeated values may be suppressed when a
smaller Artin orbit is displayed.  The word ``packet'' carries no
algebraicity assumption.  Stark's rank-one conjecture predicts
algebraic units with the prescribed Artin action, but that conjecture
is not known in this generality.  The purpose of this paper is to
identify explicit packets only where a proved theorem already supplies
algebraicity and where the remaining analytic-to-algebraic comparison
can be certified.

The word \emph{unconditional} has this precise meaning throughout:
each displayed identity follows from a published proved theorem under
hypotheses checked exactly in the case bundle.  It does not mean that
the general real-quadratic rank-one Stark conjecture is proved.
Shintani's theorem is used only in its index-two range
\cite{Shintani1978}, and the quadratic cases use the analytic
class-number formula.  Numerical PARI recognition, when retained as a
cross-check, is outside the proof chain.

\subsection*{Terminology and quantitative conventions}

Fourier inversion of (1) involves precisely the ray characters
\(\chi\) for which \(\chi(R)=-1\).  The \emph{support order} is the
maximum order of such a character.  Labels such as
\texttt{RQ-000190} are stable locators in the accompanying corpus;
they are not mathematical notation, and every use below also states
the field and modulus.

For Engine B, a \emph{safe exponent} is an explicitly computed
integer \(m\) for which Shintani's theorem proves that every
\(X_A^m\) is algebraic in the asserted ray field.  If the analytic
and algebraic logarithms differ by balls of radius at most \(\delta\),
the height comparison bounds the quotient by \(m\delta\).  The
reported \emph{margin} is
\[
 \mathcal M=\frac{V_{\min}}{m\delta},
\tag{2}
\]
where \(V_{\min}\) is the minimum applicable Voutier lower bound over
the certified degree window.  Thus \(\mathcal M>100\) is the
pre-registered acceptance condition.  A \emph{cone evaluator} is the
finite Shintani--Yamamoto decomposition of a partial zeta derivative
into double-sine or Mellin-integral terms; all such terms below are
evaluated as Arb balls.

\subsection*{How the results fit together}

There are two theorem mechanisms.  Engine A is a uniform closed-form
formula for quadratic Fourier support.  Engine B reaches higher
character orders when Shintani's maximal absolutely abelian subfield
has index two: Proposition~\ref{prop:transfer} translates the
algebraicity theorem and Lemma~\ref{lem:height} turns certified
archimedean comparisons into exact identities.  The order-six cases
on either side of ramification isolate the index/wildness condition,
rather than character order, as the relevant obstruction.
Lemmas~\ref{lem:noD} and \ref{lem:parity} explain two further
structural boundaries.

\subsection*{Main theorems}

\begin{theorem}[Uniform quadratic-support theorem]\label{thm:A}
For every real quadratic \(K\), every one-place modulus
\(\mathfrak m=\mathfrak f\infty_2\), and every packet whose
differenced Fourier support consists only of quadratic characters,
each \(X_A\) is an explicit rational power of a product of relative
units.  Corollary~\ref{cor:Aclosed} gives one closed product formula
whose exponents are expressed entirely in terms of exact class
numbers, roots of unity, imprimitive Euler factors, and
relative-regulator indices.  No numerical approximation is required
to state or prove the formula.
\end{theorem}

\begin{theorem}[Selected higher-order packets]\label{thm:main}
Use the integral basis \(1,\omega\), where
\(\omega=(1+\sqrt D)/2\) for \(D\equiv1\pmod4\) and
\(\omega=\sqrt D\) otherwise, and let the displayed HNF matrix denote
the finite ideal \(\mathfrak f\) in that basis.  For every row below, the packet
\(\mathcal X_{\mathfrak f\infty_2}\) is the Artin-labeled set of
positive roots of the relative polynomial printed in
Section~\ref{sec:packets}.
\begin{center}\scriptsize
\begin{tabular}{@{}lllcrr@{}}
\toprule
case & \(K\) & \(\operatorname{HNF}(\mathfrak f)\) & support
 & \(m\) & margin\\
\midrule
RQ-000190 & \(\Q(\sqrt7)\) &
\(\left[\begin{smallmatrix}7&0\\0&1\end{smallmatrix}\right]\)&
\(2,6\)&4032&\(>5688\)\\
RQ-000419 & \(\Q(\sqrt{14})\) &
\(\left[\begin{smallmatrix}7&0\\0&1\end{smallmatrix}\right]\)&
\(2,6\)&4032&\(>7315\)\\
RQ-000108 & \(\Q(\sqrt5)\) &
\(\left[\begin{smallmatrix}15&6\\0&3\end{smallmatrix}\right]\)&
\(4\)&2880&\(>2460\)\\
RQ-000021 & \(\Q(\sqrt2)\) &
\(\left[\begin{smallmatrix}7&0\\0&7\end{smallmatrix}\right]\)&
\(2,6\)&2016&\(>4261\)\\
RQ-002057 & \(\Q(\sqrt{57})\) &
\(\left[\begin{smallmatrix}9&3\\0&3\end{smallmatrix}\right]\)&
\(2,6\)&2592&\(>748\)\\
RQ-002955 & \(\Q(\sqrt{77})\) &
\(\left[\begin{smallmatrix}7&3\\0&1\end{smallmatrix}\right]\)&
\(2,6\)&4032&\(>5151\)\\
RQ-001107 & \(\Q(\sqrt{33})\) &
\(\left[\begin{smallmatrix}11&5\\0&1\end{smallmatrix}\right]\)&
\(2,10\)&15840&\(>5817\)\\
\bottomrule
\end{tabular}
\end{center}
Here \(m\) is the safe exponent of
Proposition~\ref{prop:transfer}, and ``margin'' has the meaning in
(2).  Each row includes exact ray-field identification, unconditional
class and unit groups, Sturm isolation, all-embeddings Arb bounds, and
exact Artin labels.
\end{theorem}

\paragraph{Historical scope.}
A dated literature search found no earlier unconditional oriented
examples with support order six or ten, so these are apparently the
first examples in those order classes.  This historical observation
is not part of either theorem and is not asserted as a universal
priority claim.\footnote{The
search record names the sources, versions, hashes, and date.  Kopp's
arXiv:2411.06763 was checked again on 30 July 2026; the current version
remains v3, dated 3 May 2025, and its algebraicity applications are
conditional.}

Theorem~\ref{thm:main} is deliberately a selected-results theorem,
not a census-completeness statement.  Counts, conductor trends, and a
complete frontier taxonomy require the separately audited census and
will appear elsewhere.  This separation prevents provisional
classification data from entering a theorem proof.  The prior
numerical and conditional computations of
Dummit--Sands--Tangedal and Cohen--Roblot remain prior work rather
than claims displaced here
\cite{DST1997,DST2003,CohenRoblot2000,Roblot2013}.
More precisely, the 1997 Dummit--Sands--Tangedal computations concern
totally real cubic bases and are numerical; their 2003 exponent-two
theorems overlap the theorem class of Engine A, which is consequently
not advertised as a historical first.  Cohen--Roblot construct
Hilbert class fields conditionally from Stark units and verify the
resulting fields independently, whereas every promoted object here is
a one-place ray field with nontrivial finite modulus.  Thus these are
mathematically adjacent precedents, not identical class-field
objects.  The order-six and order-ten historical statements refer
only to the dated search described above.

\section{The quadratic-support theorem}

Let \(G=\Cl_{\mathfrak m}(K)\), and suppose every character
\(\chi\) with \(\chi(R)=-1\) has order two.  Write
\(\chi^\circ\) for the associated primitive Hecke character and
\(L_\chi/K\) for the quadratic extension it cuts out.  The selected
real place becomes complex in \(L_\chi\), while the other real place
splits, so \(L_\chi\) has signature \((2,1)\).

Let \(\epsilon_K\) generate \(E_K/\mu(K)\).  Modulo torsion, use the
logarithmic embedding
\[
 \lambda_L(u)=
 \bigl(\log|\tau_1u|,\log|\tau_2u|,
       2\log|\tau_3u|\bigr),
\tag{3}
\]
where \(\tau_1,\tau_2\) are real and \(\tau_3\) is one of the complex
embeddings.  Let \(u_\chi\) generate the primitive rank-one kernel of
\[
 N_{L_\chi/K}:E_{L_\chi}/\mu(L_\chi)
                  \longrightarrow E_K/\mu(K),
\]
and orient it by \(\log|\tau_1u_\chi|>0\).  In any exact basis of
\(E_{L_\chi}/\mu(L_\chi)\), write \(v_K\) and \(v_\chi\) for the
coordinate columns of \(\epsilon_K\) and \(u_\chi\), and define
\[
 I_\chi=
 [E_{L_\chi}/\mu(L_\chi):
   \langle E_K/\mu(K),u_\chi\rangle]
 =|\det(v_K,v_\chi)|.
\tag{4}
\]

\begin{proposition}[Explicit quadratic formula]\label{prop:A}
With \(E_\chi\) the exact product of imprimitive Euler factors at
zero,
\[
 L'_{\mathfrak m}(0,\chi)=
 E_\chi\frac{h_{L_\chi}}{h_K}
 \frac{w_K}{w_{L_\chi}}\frac{2}{I_\chi}
 \log|\tau_1u_\chi|.
\tag{5}
\]
Moreover
\[
 Z'_{\mathfrak m}(0,A)=
 \frac{2}{|G|}
 \sum_{\chi(R)=-1}\chi(A)^{-1}
 L'_{\mathfrak m}(0,\chi).
\tag{6}
\]
Consequently a common denominator in (5)--(6) gives an explicit
power identity for every \(X_A\).
\end{proposition}

\begin{proof}
For a primitive quadratic \(\chi^\circ\),
\(\zeta_{L_\chi}(s)=\zeta_K(s)L(s,\chi^\circ)\).
The analytic class-number formula at zero, with the same regulator
normalization as (3), is
\[
 \zeta_F(s)=
 -\frac{h_FR_F}{w_F}s^{r_1(F)+r_2(F)-1}
 +O\!\left(s^{r_1(F)+r_2(F)}\right).
\]
Since the unit ranks of \(K\) and \(L_\chi\) are one and two,
respectively, division gives
\[
 L'(0,\chi^\circ)=
 \frac{h_{L_\chi}}{h_K}\frac{w_K}{w_{L_\chi}}
 \frac{R_{L_\chi}}{R_K}.
\tag{7}
\]

We next prove the regulator identity without suppressing the mixed
signature.  Put \(a=\log|\tau_1\epsilon_K|\), where
\(\tau_1,\tau_2\) lie over the split real place of \(K\).  The
product formula gives
\[
 \lambda_L(\epsilon_K)=(a,a,-2a).
\]
For the relative unit \(u_\chi\), exact norm one modulo torsion gives
\[
 \lambda_L(u_\chi)=(b,-b,0),\qquad
 b=\log|\tau_1u_\chi|>0.
\]
The determinant of the first two coordinates is \(2|a|b\).
Since \(R_K=|a|\), the regulator of
\(\langle E_K,u_\chi\rangle\) is \(2R_K\log|\tau_1u_\chi|\).
The same subgroup has index \(I_\chi\) in the full free unit group,
so its regulator is \(I_\chi R_{L_\chi}\).  Therefore
\[
 \frac{R_{L_\chi}}{R_K}
 =\frac{2}{I_\chi}\log|\tau_1u_\chi|.
\tag{8}
\]

If \(\mathfrak m/\mathfrak f_{\chi^\circ}\) is the imprimitive part,
then
\[
 L_{\mathfrak m}(s,\chi)
 =L(s,\chi^\circ)
  \prod_{\mathfrak p\mid
    \mathfrak m/\mathfrak f_{\chi^\circ}}
  \left(1-\chi^\circ(\mathfrak p)N\mathfrak p^{-s}\right).
\]
At \(s=0\) the product is the exact integer \(E_\chi\).  If one
factor vanishes, both this product and \(L(s,\chi^\circ)\) vanish at
zero, so \(L'_{\mathfrak m}(0,\chi)=0\); otherwise combining (7) and
(8) proves (5).

Finally, character inversion gives
\[
 \zeta_{\mathfrak m}(s,A)
 =\frac1{|G|}\sum_{\chi\in\widehat G}
       \chi(A)^{-1}L_{\mathfrak m}(s,\chi).
\]
Subtracting the formula with \(A\) replaced by \(RA\) multiplies the
\(\chi\)-term by \(1-\chi(R)^{-1}\), which is two precisely for
\(\chi(R)=-1\) and zero otherwise.  Differentiation proves (6).
Every surviving character is quadratic, hence every coefficient is
rational.  If \(q\) clears their denominators, exponentiation yields
\[
 X_A^q=\prod_{\chi(R)=-1}
        |\tau_1u_\chi|^{\,n_{A,\chi}}
 \qquad(n_{A,\chi}\in\Z),
\tag{9}
\]
after the exact factors in (5) have been absorbed into the integers
\(n_{A,\chi}\).  This proves Proposition~\ref{prop:A} and
Theorem~\ref{thm:A}.
\end{proof}

\begin{corollary}[Closed packet formula]\label{cor:Aclosed}
Under the hypotheses of Theorem~\ref{thm:A}, put
\[
 c_{A,\chi}
 =\frac{4}{|G|}\,\chi(A)^{-1}E_\chi
   \frac{h_{L_\chi}}{h_K}
   \frac{w_K}{w_{L_\chi}}
   \frac1{I_\chi}.
\]
Then the positive real invariant at the selected embedding is
\[
 \boxed{\qquad
 X_A=\prod_{\chi(R)=-1}
       |\tau_1u_\chi|^{\,c_{A,\chi}}.
 \qquad}
\]
Every exponent \(c_{A,\chi}\) is rational and is computed by exact
finite arithmetic.
\end{corollary}

\begin{remark}
The word ``uniform'' refers to this single character product, not to
one field-independent named unit.  For fixed
\((K,\mathfrak m,A)\), the supported characters form a finite set,
and substituting their exact class numbers, Euler factors, unit
kernels, and determinant indices reduces the corollary immediately
to a concrete product of explicitly computed relative units.
\end{remark}

The theorem is uniform; its proof does not enumerate the census.
There are 1,560 nontrivial eligible occurrences in the frozen range,
involving 2,232 supported character occurrences and 912 distinct
quartic fields.  These figures describe a verification queue, not a
completeness claim here.  Each application still certifies its
primitive conductor, field, class numbers, Euler factors, exact norm
map, primitive kernel, determinant, and orientation.  The
dimension-four and dimension-eight anchors independently replay (5),
including relative index \(I_\chi=2\) in both dimension-eight cases.

\section{Shintani transfer and height rigidity}

Let \(H=H_{\mathfrak m}\) be the one-place ray class field of
\(\mathfrak m=\mathfrak f\infty_2\), and let \(\mathcal N\) be its
normal closure over \(\Q\).  Engine B first checks Shintani's
conditions (0-3), (0-6), and (0-9), including
\[
 [H:H\cap\Q^{\rm ab}]=2.
\tag{10}
\]
The maximal absolutely abelian subfield is the fixed field of
\([\Gal(\mathcal N/\Q),\Gal(\mathcal N/\Q)]\).  The imaginary
quadratic field and transfer conductor below are then reconstructed a
second time from their intrinsic ray class fields.

More precisely, an Engine-B transfer route consists of an imaginary
quadratic subfield \(k\subset\mathcal N\), a finite modulus
\(\mathfrak c\) of \(k\), and a \(k\)-isomorphism
\[
 \iota:H_{\mathfrak c}(k)\xrightarrow{\;\sim\;}\mathcal N.
\]
Here \(\mathcal N/k\) is abelian,
\(\mathfrak c=\operatorname{cond}(\mathcal N/k)\), and
\(H_{\mathfrak c}(k)\) denotes the corresponding ray class field.
Thus ``transfer conductor'' below means this Artin conductor, not a
conductor guessed from a discriminant.  In every selected case it is
computed by \texttt{rnfconductor} and the isomorphism is independently
constructed by \texttt{bnrclassfield}; the two routes printed for a
case must produce the same \(\mathcal N\).

\begin{proposition}[Shintani transfer in one-place notation]
\label{prop:transfer}
Let \(C=\Cl_{\mathfrak f\infty_1\infty_2}(K)\), and suppose exact
fiber identities identify Shintani's invariant
\(X_{\mathfrak f}(c,G)\) with \(X_A\) in (1).  Assume Shintani's
conditions (0-3), (0-6), and (0-9), and let \(k\) and
\(\mathfrak c\) be respectively the imaginary quadratic field and
finite transfer conductor occurring in his proof.

For \(\mathfrak d\mid\mathfrak c\), let
\[
 h(\mathfrak d)=|\Cl_{\mathfrak d}(k)|,\qquad
 w(\mathfrak d)
 =|\{\zeta\in\mu(k):\zeta\equiv1\pmod{\mathfrak d}\}|,
\]
put
\[
 n(\mathfrak d)
 =w(\mathfrak d)\frac{h(\mathfrak c)}{h(\mathfrak d)},
\tag{11}
\]
and, for \(\mathfrak d\ne\mathcal O_k\), let
\(f_{\mathfrak d}\) be the positive generator of
\(\mathfrak d\cap\Z\).  If \(r\) clears the explicitly computed
real-side distribution and induction denominators, then
\[
\begin{split}
 m_{\mathcal O_k}&=12h_k n(\mathcal O_k),\\
 m_{\mathfrak d}&=12f_{\mathfrak d}n(\mathfrak d)
 \quad(\mathfrak d\ne\mathcal O_k),\\
 m&=\operatorname{lcm}\bigl(r,\,
       (m_{\mathfrak d})_{\mathfrak d\mid\mathfrak c}\bigr)
\end{split}
\tag{12}
\]
has the following property:
\[
 X_A^m\in\mathcal O_H^\times,\qquad
 (X_A^m)^{\operatorname{Art}(B)}=X_{AB}^m.
\tag{13}
\]
At every embedding of \(H\) above the nonsplit real place,
\(|X_A^m|=1\).
\end{proposition}

\begin{proof}
Shintani's Theorem~1 identifies the logarithm of
\(X_{\mathfrak f}(c,G)\) with the full-ray zeta differences over the
two fibers of \(A\) and \(RA\); the assumed exact fiber identity
therefore gives \(X_A\).  Under (0-3), (0-6), and (0-9),
Shintani's Theorem~2 supplies the index-two algebraicity construction
\cite{Shintani1978}.

In Shintani's equation (8), the distribution index is exactly (11).
For \(\mathfrak d\ne\mathcal O_k\), equations (4) and (9) contain
\(|\varphi_{\mathfrak d}|^{1/(6f_{\mathfrak d})}\) and then the
power \(1/n(\mathfrak d)\).  Replacing an absolute value by the
product with its conjugate contributes a factor two, so
\(12f_{\mathfrak d}n(\mathfrak d)\) clears the row.  At the trivial
divisor, equation (6) replaces \(f_{\mathfrak d}\) by \(h_k\).
Proposition~4 is an integral product of ratios of these
\(W_{\mathfrak c}\)-invariants, so it introduces no further rational
power.  The separately computed integer \(r\) clears the real-side
distribution or induction step.  Raising to (12) therefore clears
every exponent in the construction.  Shintani's Proposition~5 gives
the unit assertion, the Artin covariance, and modulus one at the
nonsplit real embeddings, proving (13).
\end{proof}

For the headline case \(K=\Q(\sqrt7)\),
\(\mathfrak f=\left[\begin{smallmatrix}7&0\\0&1\end{smallmatrix}\right]\),
the short transfer route has \(k=\Q(i)\) and
\(\mathfrak c=(7)\).  The complete divisor table is
\[
\begin{array}{c|r|r|r|r|r}
\mathfrak d&h(\mathfrak d)&w(\mathfrak d)&n(\mathfrak d)
 &f_{\mathfrak d}&m_{\mathfrak d}\\ \hline
\mathcal O_k&1&4&48&1&576\\
(7)&12&1&1&7&84
\end{array}
\tag{14}
\]
and the real-side indices are \(6\) and \(1\).  Hence
\(\operatorname{lcm}(576,84,6)=4032\).  This is a complete instance
of Proposition~\ref{prop:transfer}; the longer
\(\Q(\sqrt{-7})\) reconstruction gives the same normal closure and
safe power independently.

\begin{lemma}[All-embeddings height rigidity]\label{lem:height}
Let \(H\) be a number field of degree \(D\).  Suppose
\((X_A^m)_A\subset\mathcal O_H^\times\) is an algebraic family with
Artin covariance, and let
\((\alpha_A)_A\subset\mathcal O_H^\times\) be an exactly
Artin-labeled positive candidate family.  For each archimedean place
\(v\) of \(H\), suppose either
\[
 \left|\log|X_A|_v-\log|\alpha_A|_v\right|\le\epsilon_v
\tag{15}
\]
or both powered absolute values are proved to be one.  Put
\[
 B_A=\frac{m}{D}\sum_{v\mid\infty}n_v\epsilon_v,
\qquad n_v=[H_v:\mathbb R],
\tag{16}
\]
with \(\epsilon_v=0\) in the second case.  If
\[
 B_A<
 \min\left\{\frac12\log\frac{1+\sqrt5}{2},\
 \min_{3\le d\le D}
 \frac1{4d}\left(\frac{\log\log d}{\log d}\right)^3
 \right\},
\tag{17}
\]
then \(X_A=\alpha_A\).
\end{lemma}

\begin{proof}
The quotient
\(\eta_A=X_A^m/\alpha_A^m\) is a unit of \(H\).  Its finite-place
contribution to the Weil height is zero.  Equations (15)--(16) and
the definition of logarithmic height give
\[
 h(\eta_A)
 =\frac1D\sum_{v\mid\infty}
 n_v\max(0,\log|\eta_A|_v)
 \le B_A.
\]
If \(d=[\Q(\eta_A):\Q]\ge3\) and \(\eta_A\) is not torsion,
Voutier's explicit lower bound contradicts (17)
\cite{Voutier1996}.  A non-torsion quadratic unit has height at least
\(\frac12\log((1+\sqrt5)/2)\), and a rational unit is \(\pm1\).
Thus \(\eta_A\) is a root of unity.  At the distinguished real
embedding both \(X_A\) and \(\alpha_A\) are positive, so
\(\eta_A=1\); taking the unique positive \(m\)-th root proves the
claim.
\end{proof}

For every row of Theorem~\ref{thm:main}, the embeddings over the
split real place are all represented by the printed packet roots and
the cone enclosures.  At the remaining archimedean places,
Proposition~\ref{prop:transfer} gives modulus one for \(X_A^m\), and
the reciprocal candidate polynomial is checked exactly, after the
substitution \(T=X+X^{-1}\), to have its remaining roots on the unit
circle.  Hence these places contribute \(\epsilon_v=0\) in (16).

The algebraic candidates are certified without a GRH assumption:
each class and unit computation uses \texttt{bnfinit(P,1)}, followed
by the full call \texttt{bnfcertify(bnf)} returning \(1\); the
quotient-only flag is not used.  Exact \texttt{rnfconductor} and
\texttt{bnrclassfield} round trips identify the relative ray fields,
Sturm sequences isolate the positive roots, and split-prime
Frobenius actions label them.

\paragraph{Effective but not uniformly cheap.}
The required raw radius is \(O(V_{\min}/m)\).  Arb precision therefore
grows with \(\log m\), but conductor size, packet degree, and cone
count also grow, and no uniform complexity claim is made.  The
order-ten case below has \(m=15840\) and remains comfortable; cases
with much larger exponents can exceed the fixed per-case resource cap
and belong to the frontier.  This is a stated limit of Engine B, not
evidence against its completed cases.

\section{The order-six and order-ten packets}\label{sec:packets}

\subsection{\texorpdfstring{\(\Q(\sqrt7)\)}{Q(sqrt(7))} at the
ramified prime above seven}

For \(K=\Q(\sqrt7)\) and
\(\mathfrak m=\mathfrak p_7\infty_2\), the one-place and two-place ray
groups are \(C_6\) and \(C_6\times C_2\).  The support orders are two
and six and (10) holds.  The commutator-fixed calculation produces
\(\Q(i)\) and \(\Q(\sqrt{-7})\), and their intrinsic conductor
reconstructions give the same degree-$24$ normal closure.  The
\(\Q(i)\) route has two divisor rows and safe exponent \(4032\).

The packet polynomial over \(K\) is
\[
\begin{aligned}
P_7(X)={}&X^6-(6+2\sqrt7)X^5+(22+8\sqrt7)X^4\\
 &-(34+13\sqrt7)X^3+(22+8\sqrt7)X^2\\
 &-(6+2\sqrt7)X+1.
\end{aligned}
\tag{19}
\]
It is irreducible over \(K\), defines the one-place ray field, has six
positive Sturm-isolated roots, and is Artin-labeled by the split
primes \(19\) and \(31\).

For the class represented by \((3)^r\), an exact Yamamoto domain has
generators
\[
 \lambda_r=(7-2\sqrt7)/3^r,\qquad
 \lambda_r(8+3\sqrt7).
\]
Its determinant is nine and its half-open parallelotope contains
exactly nine lattice points.  The resulting 27 independent
double-sine evaluations are enclosed without a PARI \(L\)-value.
At exponent \(4032\), the powered height upper bound is
\(9.1903\times10^{-9}\).  The minimum Voutier lower bound in degrees
three through 24 is \(5.2279\times10^{-5}\), so the margin exceeds
5688.  The degree-one and degree-two fallbacks finish the proof of
(19).

This is the headline order-six instance.  The same character order is
the difficult component at the separate \(\Q(\sqrt{21})\) wall, but
here the index-two hypothesis holds and the packet closes.  Thus the
character order itself is not the obstruction.

\subsection{Replications and the
\texorpdfstring{ramified-prime-\(3\)}{ramified-prime-3} control}

For \(\Q(\sqrt{14})\) at \(\mathfrak p_7\), both
\(\Q(\sqrt{-7})\) and \(\Q(\sqrt{-2})\) reconstruct the degree-$24$
normal closure.  The latter route has clearing exponents 576 and 84,
with least common multiple 4032.  The packet is
\[
\begin{aligned}
P_{14}(X)={}&X^6-(13+4\sqrt{14})X^5
 +(85+22\sqrt{14})X^4\\
&-(139+38\sqrt{14})X^3
 +(85+22\sqrt{14})X^2\\
&-(13+4\sqrt{14})X+1.
\end{aligned}
\tag{20}
\]
Split primes 11 and 103 label the Artin generators.  The 20-point
cone evaluator gives powered height at most
\(7.1466\times10^{-9}\), hence margin greater than 7315.

The strongest boundary control is RQ-002057 over
\(\Q(\sqrt{57})\), with modulus of norm \(27\).  The conductor is a
power of the ramified prime 3 and the support still has order six.
The bases \(\Q(\sqrt{-3})\) and \(\Q(\sqrt{-19})\) independently
give the same degree-$24$ closure.  On the shorter route the divisor
exponents are \(864,324,108\), hence \(m=2592\).
With \(y=(1+\sqrt{57})/2\), so \(y^2-y-14=0\), its packet is
\[
\begin{aligned}
P_{57}(X)={}&X^6-(30+9y)X^5+(582+177y)X^4\\
&-(2529+772y)X^3+(582+177y)X^2\\
&-(30+9y)X+1.
\end{aligned}
\tag{20a}
\]
The fundamental positive unit has order three modulo the finite
modulus, so a ray-class domain is the union of three adjacent cones,
with 240 exact affine points per class.  The powered height is at most
\(6.982\times10^{-8}\), giving margin greater than 748.  Thus
order-six support and ramified-prime-$3$ structure are both
reachable; neither alone explains the separate
\(\Q(\sqrt{21})\) boundary.

\subsection{Three further Shintani closures}

The three remaining rows of Theorem~\ref{thm:main} use exactly the
Engine-B proof above; we print their routes and algebraic candidates
so that no theorem row is supported only by a corpus identifier.

For RQ-000108, \(K=\Q(\sqrt5)\) and
\(\Norm\mathfrak f=45\).  The two imaginary bases are
\(\Q(\sqrt{-3})\) and \(\Q(\sqrt{-15})\), and both reconstruct the
same degree-$16$ normal closure.  With
\(y=(1+\sqrt5)/2\), the relative packet polynomial is
\[
 X^4-(1+6y)X^3+(9+9y)X^2-(1+6y)X+1.
\tag{21}
\]
It is the certified one-place ray field over \(K\).
At \(m=2880\), the powered height is at most
\(2.12465\times10^{-8}\), so the Voutier margin exceeds 2460.

For RQ-000021, \(K=\Q(\sqrt2)\) and
\(\mathfrak f=(7)\).  The bases \(\Q(\sqrt{-7})\) and
\(\Q(\sqrt{-14})\) reconstruct the same degree-$24$ closure.  The
packet polynomial is
\[
\begin{aligned}
X^6&-(10+6\sqrt2)X^5+(58+40\sqrt2)X^4\\
&-(129+90\sqrt2)X^3+(58+40\sqrt2)X^2\\
&-(10+6\sqrt2)X+1.
\end{aligned}
\tag{22}
\]
The cone boundary meets a split-prime face; the canonical value is
the average of the upper and lower half-open conventions.  Both are
enclosed independently.  At \(m=2016\), the powered height is at most
\(1.22671\times10^{-8}\), giving margin greater than 4261.

For RQ-002955, \(K=\Q(\sqrt{77})\) and
\(\mathfrak f=\mathfrak p_7\) has HNF
\(\left[\begin{smallmatrix}7&3\\0&1\end{smallmatrix}\right]\).
The bases \(\Q(\sqrt{-7})\) and \(\Q(\sqrt{-11})\) give the same
degree-$24$ closure.  Put \(y=(1+\sqrt{77})/2\), so
\(y^2-y-19=0\).  The packet is
\[
\begin{aligned}
X^6&-(15+3y)X^5+(107+26y)X^4\\
&-(217+54y)X^3+(107+26y)X^2-(15+3y)X+1.
\end{aligned}
\tag{23}
\]
At \(m=4032\), the powered height is at most
\(1.01485\times10^{-8}\), and the margin exceeds 5151.  In all three
cases the exact ray-field identification, Sturm isolation, Artin
labels, and the degree-one and degree-two fallbacks are included in
the cited case records.

\subsection{Order ten}

For RQ-001107 over \(\Q(\sqrt{33})\) and
\(\mathfrak p_{11}\), put \(y=(1+\sqrt{33})/2>0\), so
\(y^2-y-8=0\).  The packet is
\[
\begin{aligned}
P_{33}(X)={}&X^{10}-(8+4y)X^9+(74+30y)X^8
 -(294+125y)X^7\\
&+(669+281y)X^6-(871+368y)X^5
 +(669+281y)X^4\\
&-(294+125y)X^3+(74+30y)X^2-(8+4y)X+1.
\end{aligned}
\tag{24}
\]
Its relative degree is ten, so the certified ray field \(H/K\) has
absolute degree \(20\).  The normal closure used to compute the
Frobenius action has degree \(40\), but the comparison quotient
\(\eta_A=X_A^m/\alpha_A^m\) lies in \(H\), since both powered terms
do.  Therefore Lemma~\ref{lem:height} needs only \(d\le20\).  The
archived computation conservatively evaluated every degree through
80; on either window the minimum Voutier bound is
\[
 5.22795332222684\ldots\times10^{-5}.
\]
At \(m=15840\), the factor-100 raw-error target is
\(3.30047558221391\ldots\times10^{-11}\); the achieved bound is
\(5.6731\times10^{-13}\), so the certified margin exceeds 5817.
The degree-one and degree-two fallbacks are recorded separately.
This support order is the first one in the present proved inventory
divisible by the prime \(5\); its historical priority status is the
survey-bounded status stated after Theorem~\ref{thm:main}.

\section{Structural boundary lemmas}

\begin{lemma}[Absolutely abelian ray fields]\label{lem:noD}
Let \(H/K\) be the one-place ray class field of a real quadratic
field.  If \(H/\Q\) is Galois and abelian, then the sign class
\(R\in\Gal(H/K)\) is trivial and the differenced Fourier support in
(1) is empty.
\end{lemma}

\begin{proof}
If the one-place field \(H/\Q\) were abelian, the complex
conjugations above the two real places of \(K\) would be conjugate:
any lift of the nontrivial element of \(\Gal(K/\Q)\) interchanges the
places.  In an abelian group conjugate elements are equal.  The
inertia element at the omitted real place is trivial in the one-place
ray field, whereas the inertia element at the retained place is
\(R\).  Hence \(R=1\).  It follows immediately that no character
satisfies \(\chi(R)=-1\), so the differenced Fourier support is empty.
\end{proof}

Engine A is not an exception: Proposition~\ref{prop:A} reduces
individual quadratic characters through (5)--(6), but it does not
assert that the entire mixed-signature one-place ray field is
absolutely abelian.

\begin{lemma}[Index parity]\label{lem:parity}
Let \(H/K\) be a one-place ray field and let \(N/\Q\) be its genuine
normal closure.  If the sign class \(R\) is nontrivial, then
\[
 2\mid[H:H\cap\Q^{\rm ab}].
\tag{32}
\]
Equivalently, an odd Shintani index can occur only when the Fourier
support of (1) is empty.
\end{lemma}

\begin{proof}
Put \(G=\Gal(N/\Q)\), \(A=\Gal(N/K)\), and
\(B=\Gal(N/H)\).  Let \(c_1,c_2\in A\) be the real-place inertia
involutions.  A lift of the nontrivial automorphism of \(K/\Q\)
interchanges them, so they are conjugate in \(G\).  Hence
\(c_1c_2^{-1}\in[G,G]\).  Since \(G/A\simeq C_2\), one has
\([G,G]\subseteq A\).  Moreover \(B\) is normal in \(A\), so
\([G,G]B\) is a subgroup of \(A\).  In \(A/B=\Gal(H/K)\), the
involution at the
omitted place is trivial and the one at the retained place is \(R\).
Thus the image of \([G,G]\) in \(A/B\) contains \(R\).  Its order is
the degree of \(H\) over its fixed field.  That fixed field is
\[
 H\cap N^{[G,G]}=H\cap\Q^{\rm ab},
\]
because \(N^{[G,G]}\) is the maximal abelian subextension of
\(N/\Q\).  Hence the image has order
\([H:H\cap\Q^{\rm ab}]\), which is even if \(R\ne1\).
Finally, characters of a finite abelian group separate a nonidentity
element of order two, so \(R=1\) is equivalent to empty differenced
support.
\end{proof}

The use of the genuine normal closure in Lemma~\ref{lem:parity} is
essential.  Conjugation coordinates at an unstable modulus do not
define \([G,G]\).  As an independent consistency check, a
pre-registered replay over all 8,200 genuine normal-closure
reconstructions found 446 odd indices greater than one.  Every one
had trivial sign class and empty differenced Fourier support; there
were no exceptions.

\appendix

\section{Certificates and provenance}

The exact field computations use PARI/GP 2.15.4; analytic
certificates use Python 3.12.3 and python-flint 0.9.0.  The principal
replays were run under Linux/x86-64 on an AMD EPYC 9354P virtual CPU.
All analytic conclusions are Arb balls, never point estimates.
AI-assisted systems proposed code, text, and adversarial checks but
were not mathematical authorities: every accepted computation was
rerun by the pinned exact or Arb pipeline and reviewed by the author.
The disclosure immediately before the bibliography gives the tools
and purposes.

The principal machine-readable records are
\begin{center}\small
\begin{tabular}{@{}p{0.29\linewidth}p{0.65\linewidth}@{}}
\toprule
result & record\\
\midrule
first order six & \path{data/q7-p7-case-v1.json}\\
second order six & \path{data/q14-p7-case-v1.json}\\
RQ-000108 & \path{data/rq000108-case-v1.json}\\
RQ-000021 & \path{data/rq000021-case-v1.json}\\
ramified-prime-$3$ control & \path{data/q57-norm27-case-v1.json}\\
RQ-002955 & \path{data/rq002955-case-v1.json}\\
order ten & \path{data/q33-p11-order10-case-v1.json}\\
all archimedean places &
\path{artifacts/engine-b-archimedean-place-audit-v1.json}\\
quadratic theorem & \path{data/engine-a-uniform-theorem-v1.json}\\
index-parity theorem & \path{artifacts/results-paper-index-parity-lemma-v1.json}\\
parity audit & \path{artifacts/results-paper-odd-index-parity-audit-v1.json}\\
literature perimeter and Kopp watch &
\path{data/literature-perimeter-v1.json};
\path{artifacts/kopp-arxiv-watch-2026-07-30.json}\\
\bottomrule
\end{tabular}
\end{center}

\subsection*{Artin-label replay}

For each cyclic ray group below, \(r\) denotes the power of the
generator in the pinned PARI ray coordinates.  The entry following
\(r\) is a rational Sturm interval containing exactly one positive
root of the absolute packet polynomial.  The cone evaluation for
class \(r\), Shintani's Artin covariance (13), and the exact
ray-field isomorphism identify that root with \(X_r\).
\begin{center}\scriptsize
\begin{tabular}{@{}p{0.14\linewidth}p{0.80\linewidth}@{}}
\toprule
case & \(r:\) isolating interval\\
\midrule
RQ-000190 &
\(0:(5.325,5.326),\ 1:(2.713,2.714),\
2:(2.251,2.252),\ 3:(.187,.188),\
4:(.368,.369),\ 5:(.444,.445)\)\\
RQ-000419 &
\(0:(20.431,20.432),\ 1:(.189,.190),\
2:(1.090,1.091),\ 3:(.048,.049),\
4:(5.289,5.290),\ 5:(.916,.917)\)\\
RQ-000108 &
\(0:(7.892,7.893),\ 1:(2.242,2.243),\
2:(.126,.127),\ 3:(.445,.446)\)\\
RQ-000021 &
\(0:(6.109,6.110),\ 1:(7.496,7.497),\
2:(4.351,4.352),\ 3:(.163,.164),\
4:(.133,.134),\ 5:(.229,.230)\)\\
RQ-002057 &
\(0:(5.709,5.710),\ 1:(.028,.029),\
2:(.035,.036),\ 3:(.175,.176),\
4:(34.732,34.733),\ 5:(27.792,27.793)\)\\
RQ-002955 &
\(0:(18.232,18.233),\ 1:(2.180,2.181),\
2:(8.620,8.621),\ 3:(.054,.055),\
4:(.458,.459),\ 5:(.116,.117)\)\\
RQ-001107 &
\(0:(8.645,8.646),\ 1:(.276,.277),\
2:(1.245,1.246),\ 3:(.571,.572),\
4:(.236,.237),\ 5:(.115,.116),\
6:(3.611,3.612),\ 7:(.802,.803),\
8:(1.748,1.749),\ 9:(4.234,4.235)\)\\
\bottomrule
\end{tabular}
\end{center}

Each record carries hashes of its exact and analytic dependencies.
The supplementary manifest also records the full software versions,
commands, source hashes, and failed computations; those provenance
details are not part of the mathematical argument.

\paragraph{Data availability.}
The companion archive will be deposited before submission.  It will
contain the source, exact and Arb certificates, the complete
SHA-256 manifest, and one-command replay instructions.

\section*{Declaration of generative AI and AI-assisted technologies
in the manuscript preparation process}

During preparation of this work, the author used OpenAI GPT-5.6 for
mathematical brainstorming, symbolic-computation scaffolding, code
review, and manuscript editing, and Anthropic Claude for
independent adversarial checking of statements and certificates.  The
author reviewed and verified all outputs, made all mathematical and
editorial decisions, and takes full responsibility for the content of
the article.

\begin{thebibliography}{20}
\raggedright
\bibitem{Shintani1978}
T.~Shintani, \emph{On certain ray class invariants of real quadratic
fields}, J.\ Math.\ Soc.\ Japan \textbf{30} (1978), 139--167,
\url{https://doi.org/10.2969/jmsj/03010139}.
\bibitem{Stark1980}
H.~M.~Stark, \emph{\(L\)-functions at \(s=1\). IV. First derivatives
at \(s=0\)}, Adv.\ Math.\ \textbf{35} (1980), 197--235,
\url{https://doi.org/10.1016/0001-8708(80)90049-3}.
\bibitem{Yamamoto2010}
S.~Yamamoto, \emph{On Shintani's ray class invariant for totally real
number fields}, Math.\ Ann.\ \textbf{346} (2010), 559--593,
\url{https://doi.org/10.1007/s00208-009-0405-x}.
\bibitem{Tangedal2007}
B.~A.~Tangedal, \emph{Continued fractions, special values of the
double sine function, and Stark units over real quadratic fields},
J.\ Number Theory \textbf{124} (2007), 291--313,
\url{https://doi.org/10.1016/j.jnt.2006.09.011}.
\bibitem{Voutier1996}
P.~Voutier, \emph{An effective lower bound for the height of algebraic
numbers}, Acta Arith.\ \textbf{74} (1996), 81--95,
\url{https://doi.org/10.4064/aa-74-1-81-95}.
\bibitem{DST1997}
D.~S.~Dummit, J.~W.~Sands, and B.~A.~Tangedal,
\emph{Computing Stark units for totally real cubic fields},
Math.\ Comp.\ \textbf{66} (1997), 1239--1267,
\url{https://doi.org/10.1090/S0025-5718-97-00852-1}.
\bibitem{DST2003}
D.~S.~Dummit, J.~W.~Sands, and B.~A.~Tangedal,
\emph{Stark's conjecture in multiquadratic extensions, revisited},
J.\ Th\'eorie des Nombres de Bordeaux \textbf{15} (2003), 83--97,
\url{https://doi.org/10.5802/jtnb.388}.
\bibitem{CohenRoblot2000}
H.~Cohen and X.-F.~Roblot, \emph{Computing the Hilbert class field of
real quadratic fields}, Math.\ Comp.\ \textbf{69} (2000), 1229--1244,
\url{https://doi.org/10.1090/S0025-5718-99-01111-4}.
\bibitem{Roblot2013}
X.-F.~Roblot, \emph{Index formulae for Stark units and their
solutions}, Pacific J.\ Math.\ \textbf{266} (2013), 391--422,
\url{https://doi.org/10.2140/pjm.2013.266.391}.
\bibitem{Kopp2025}
G.~Kopp, \emph{The Shintani--Faddeev modular cocycle: Stark units from
\(q\)-Pochhammer ratios}, arXiv:2411.06763.
\bibitem{DasguptaKakde2023}
S.~Dasgupta and M.~Kakde, \emph{On the Brumer--Stark conjecture},
Ann.\ of Math.\ \textbf{197} (2023), 289--388,
\url{https://doi.org/10.4007/annals.2023.197.1.5}.
\bibitem{TangedalYoung2013}
B.~A.~Tangedal and P.~T.~Young, \emph{Explicit computation of
Gross--Stark units over real quadratic fields},
J.\ Number Theory \textbf{133} (2013), 1022--1045,
\url{https://doi.org/10.1016/j.jnt.2012.04.021}.
\bibitem{PARI}
The PARI Group, \emph{PARI/GP version 2.15.4},
\url{https://pari.math.u-bordeaux.fr/}.
\bibitem{pythonflint}
F.~Johansson and contributors, \emph{python-flint 0.9.0},
\url{https://github.com/flintlib/python-flint}.
\end{thebibliography}

\end{document}
